\section{Related Work}

\subsection{Learning-Based Seismic Velocity Model Building}

Learning-based seismic inversion has developed from direct convolutional mappings between seismic measurements and subsurface velocity toward larger encoder--decoder, attention-based, and structured inverse architectures. Fully convolutional approaches demonstrated the feasibility of amortized velocity inversion \citep{yang2019deeplearning}, while InversionNet established a widely used OpenFWI benchmark \citep{wu2020inversionnet,deng2022openfwi}. Later methods explored alternative architectures and physical priors, and UPFWI introduced an unsupervised, wave-equation-coupled formulation \citep{jin2022upfwi}. These studies differ substantially in observation type and optimization objective. The present task is narrower: reconstructing a velocity model from an already available set of synthetic processed/interpreted constraints rather than from raw seismic records.

\subsection{Generative Modeling for Velocity Reconstruction}

Generative models provide an alternative to deterministic regression when the inverse mapping contains structurally complex or multimodal variation. Adversarial generators can encourage sharp outputs but introduce adversarial optimization. Diffusion and score-based models learn iterative denoising or score fields \citep{ho2020ddpm,song2021scorebased}, and latent diffusion reduces generative cost by operating in a compressed representation \citep{rombach2022ldm}. Flow Matching and rectified flow learn continuous vector fields that transport a simple base distribution to a target distribution \citep{lipman2023flowmatching,liu2023rectifiedflow}. Recent geophysical work has applied diffusion-style modeling to controllable velocity synthesis and inversion \citep{wang2024controllable}.

\subsection{Heterogeneous Geophysical Constraint Integration}

Velocity model building can use information beyond raw shot gathers. RMS velocity constrains kinematics and velocity scale \citep{dix1955velocities}; migrated images provide reflector geometry \citep{stolt1978migration}; horizons summarize interpreted interfaces; and wells provide sparse local velocity control. Multimodal learning offers general tools for alignment, fusion, and missing-modality handling \citep{baltrusaitis2019multimodal,tian2020cmc,saporta2024symile}. A distinct question is how reconstruction responsibility is assigned after such evidence is fused.

BG-RFM addresses this latter question. Existing learned estimators generally predict the full velocity field directly or condition a full-field generator globally. Here all modalities contribute to role-aware condition interfaces, but deterministic background recovery and generative residual transport are assigned different downstream responsibilities and are coupled through the same predicted-background reference. This establishes a reconstruction architecture in which the observation-dependent background prediction also supplies the reference for generative correction.
